\begin{abstract}
针对物流分拣中心逐小时到货波动下的班次设计与月度用工问题，本文基于附件中30天、720个小时的进货记录，建立“日排班--限时处理--月度人员配置”三层混合整数规划模型。模型以整数小时班次起点、班次人数、实际处理量、待处理库存及低效率工时分配为变量，通过逐小时库存守恒避免将空闲产能错误结转至未来。

对于问题一，建立5个连续8小时班次下的日最少人数模型，30天共需10791人日，峰值日需537人。对于问题二，将每名工人一个小时的处理能力由25件降为10件，并增加0--12时到货必须在16时前完成的服务水平约束；最优结果为11951人日，峰值581人，比问题一增加10.75\%。对于问题三，综合峰值日下界与总工日下界，证明月度招工人数至少为581人；进一步求解含17430个0--1变量的月度排班模型，构造出581名工人的完整出勤与班次分配表，逐人满足工作23天且连续工作不超过7天，故581人的下界可达。敏感性分析表明，处理效率下降5\%时问题二总人日增加至12576，峰值升至611人。

本文模型保持了货物流与工时流的一致性，所有结论、图表和排班明细均可由附录程序与结果文件复现，可为动态物流用工决策提供量化依据。
\keywords{物流分拣\quad 混合整数规划\quad 库存流量\quad 人员排班\quad 服务水平}
\end{abstract}
